\chapter{Conclusions}
L'idea di realizzare questo progetto è nata da una nostra esigenza pratica, ovvero quella di avere una applicazione semplice e pratica per prendere appunti durante le lezioni universitarie. Fin da subito quindi la nostra motivazione all'idea di sviluppare Easy Notes è stata molto alta. Noi stessi siamo studenti e quindi abbiamo cercato di creare un'applicazione che rispondesse in modo concreto alle nostre esigenze e a quelle dei nostri colleghi. Il corso di Applicazioni e Servizi Web è stato un'occasione per conciliare la nostra idea con le conoscenze acquisite, mettendoci alla prova con la creazione di un'applicazione web moderna e completa che includesse tutti i livelli di uno stack web basato sul paradigma Javascript Everywhere. Fin dalle prime fasi di analisi dei requisiti e progettazione delle schermate dell'applicazione, il nostro obiettivo è stato quello di creare un interfaccia gradevole, minimale e facile da usare. Le funzionalità messe a disposizione dal framework Vue.js ci hanno permesso di implementare l'interfaccia in maniera modulare ed estendibile, oltre che reattiva e centrata sull'esperienza utente. La metafora del quaderno digitale, ha contribuito a rendere l'interazione intuitiva e coerente, riducendo la curva di apprendimento per l'utente finale e introducendo un elemento di scheumorfismo che risulta simpatico e può rendere più gradevole l'esperienza di utilizzo durante le lezioni. Dal punto di vista tecnico, oltre alle sfide legate all'implementazione di tutti gli aspetti classici di un'applicazione web, abbiamo avuto l'opportunità di integrare funzionalità avanzate come le notifiche push e la generazione automatica dei riassunti tramite LLM, che hanno arricchito ulteriormente l'applicazione e il nostro bagaglio di conoscenze. Durante lo sviluppo, la suddivisione dei lavori adottata e l'approccio agile hanno permesso a ciascun membro del team di toccare con mano tutti i livelli dello stack di sviluppo, dal design dell'interfaccia utente alla gestione e modellazione dei dati nel database, oltre che il design e l'implementazione delle API RESTful per la comunicazione tra frontend e backend. In conclusione siamo contenti di dire che l'esperienza di sviluppo di Easy Notes è stata molto positiva e formativa, la nostra precedente esperienza di sviluppo come gruppo ci ha permesso di lavorare in modo efficiente e collaborativo, ogni membro ha contribuito con le proprie competenze e idee, portando alla realizzazione di un prodotto finale di cui siamo molto soddisfatti. 